%==============================================================================
% Sjabloon projectvoorstel graduaatsproject
%==============================================================================
% Je gegevens staan in ../projectgegevens.tex. De inhoud van je voorstel
% schrijf je in voorstel-inhoud.tex - dat bestand wordt later ook automatisch
% als bijlage in je rapport opgenomen.

\documentclass{hogent-article}

% Voor een voorstel in het Engels: \documentclass[english]{hogent-article}
% Dat kan enkel na toestemming van je begeleider.

\addbibresource{voorstel.bib}

%%=============================================================================
%% Projectgegevens - VUL DIT ALS EERSTE IN
%%=============================================================================
%%
%% Dit is de enige plek waar je je naam, je titel, je bedrijf en je begeleiders
%% invult. Het voorstel, het rapport en de poster halen hun gegevens hier op,
%% zodat ze nooit uit elkaar kunnen lopen.
%%
%% Enkel de tekst tussen de accolades aanpassen. Laat de commandonamen staan.

%% ---------- Jij ------------------------------------------------------------
\newcommand{\studentnaam}{Zaher Alshikh Kheder}
\newcommand{\studentemail}{zaher.alshikhkheder@student.hogent.be}

%% Wordt gebruikt als bestandsnaam van je PDF's, bijvoorbeeld
%% Familienaam-Voornaam-rapport.pdf. Geen spaties of leestekens gebruiken.
\newcommand{\bestandsnaam}{Alshikhkheder-Zaher}

%% ---------- Je project -----------------------------------------------------
\newcommand{\projecttitel}{Systems monitor}

%% Een ondertitel is optioneel. Laat leeg als je er geen hebt: {}
\newcommand{\projectondertitel}{}

%% Een of twee zinnen. Wordt gebruikt op de poster en in het voorstel.
\newcommand{\projectpitch}{
 Kritieke services van beveiligingssystemen kunnen worden stopgezet, zowel door een technische storing als door een lokale administrator op de pc. Daarom ontwikkel ik een Windows-service die deze services continu bewaakt. Voor de implementatie in .NET (Core) gebruik ik het Worker Service-sjabloon in combinatie met het NuGet-pakket Microsoft.Extensions.Hosting.WindowsServices.

}

%% ---------- Waar je het doet ------------------------------------------------
\newcommand{\bedrijfsnaam}{IQGeo}

%% Je mentor op de werkvloer (het bedrijf)
\newcommand{\mentornaam}{Jonas Van Nieuwenberg}
\newcommand{\mentoremail}{jonas.vannieuwenberg@iqgeo.com}

%% Je begeleider van de opleiding (HOGENT)
\newcommand{\begeleidernaam}{Luc Vervoort}
\newcommand{\begeleideremail}{luc.vervoort@hogent.be}

%% ---------- Je repository ---------------------------------------------------
%% De link naar je eigen git-repository. Die hoort zichtbaar in je rapport en
%% op je poster: hij is een beoordeeld onderdeel van je oplevering.
\newcommand{\repolink}{https://github.com/zaheralshikhkheder/GraduaatProject}

%% ---------- Vast voor dit opleidingsonderdeel -------------------------------
%% Deze twee hoef je normaal niet aan te passen.
\newcommand{\opleidingnaam}{Graduaat in het Programmeren}
\newcommand{\academiejaarwaarde}{2026-2027}

%% Zoekwoorden voor het voorstel en de poster (3 tot 5)
\newcommand{\projectkeywords}{WindowsService, SystemMonitor, SecuritySystems}


%% Opleiding, vak en soort opdracht
\studyprogramme{\opleidingnaam}
\course{Graduaatsproject}
\assignmenttype{Projectvoorstel}

\academicyear{\academiejaarwaarde}

\title{\projecttitel}

\author{\studentnaam}
\email{\studentemail}

%% Je mentor op de werkvloer
\supervisor[Mentor]{\mentornaam{} (\bedrijfsnaam, \href{mailto:\mentoremail}{\mentoremail})}

%% De link naar je repository
\projectrepo{\repolink}

\keywords{\projectkeywords}

\begin{document}

\begin{abstract}
  % Eén samenhangende alinea van vijf tot acht zinnen, die je hele voorstel
  % samenvat: de context bij je bedrijf, het probleem, wat je gaat maken, met
  % welke technologie, en wat het bedrijf eraan heeft. Geen inleiding, geen
  % opsomming.
Kritieke services van beveiligingssystemen kunnen worden stopgezet, zowel door een technische storing als door een lokale administrator op de pc. Daarom ontwikkel ik een Windows-service die deze services continu bewaakt.Voor de implementatie in .NET (Core) gebruik ik het Worker Service-sjabloon in combinatie met het NuGet-pakket Microsoft.Extensions.Hosting.WindowsServices.
Ik ga een achtergrondservice programmeren die de kritieke services in de gaten houdt. Wanneer een van deze kritieke services op een beheerde laptop wordt stopgezet via een lokaal administratoraccount, stuurt mijn service direct een melding via e-mail naar het IT- en securityteam. Deze e-mail vermeldt de naam van de laptop, de naam van de service en de datum en het tijdstip van het incident.
De nieuwe service wordt op zijn beurt beschermd door een tweede service, die controleert of de eerste service actief blijft.
\end{abstract}

%%=============================================================================
%% De inhoud van je projectvoorstel
%%=============================================================================
%%
%% Zeven rubrieken, samen hoogstens twee bladzijden. Houd het kort: per
%% rubriek volstaan één tot drie zinnen. Een voorstel is geen rapport - het
%% moet je begeleider in vijf minuten laten beslissen of dit een geschikt
%% graduaatsproject is.
%%
%% Stem je voorstel af met je mentor VOORDAT je het indient. Je mentor
%% bevestigt dat dit is wat het bedrijf bedoelt en dat de tien werkdagen
%% haalbaar zijn.

%---------- 1. Context en probleemstelling ------------------------------------

\section{Context en probleemstelling}%
\label{vst:context}

% Waar loop je stage, wat doet dat team, en welk concreet probleem of welke
% concrete behoefte is de aanleiding? Noem de mensen voor wie het een probleem
% is. Vermijd algemeenheden als "bedrijven hebben nood aan".

De stage vindt plaats binnen het engineeringteam van IQGeo, dat verantwoordelijk is voor softwareontwikkeling. Het project richt zich op de beveiligingssystemen die op elke beheerde laptop geïnstalleerd zijn. Deze systemen maken gebruik van gekoppelde achtergrondservices die via een lokaal administratoraccount gestopt kunnen worden. Het IT-team wil deze services daarom graag actief monitoren. 

%---------- 2. Doelstelling ---------------------------------------------------

\section{Doelstelling}%
\label{vst:doelstelling}

% Wat bestaat er op het einde van dit project dat er vandaag niet is? Formuleer
% het als een resultaat waarvan iemand kan vaststellen dat het er is, en
% vermeld waaraan je zal merken dat het geslaagd is.

Ik ga een achtergrondservice programmeren die de kritieke services in de gaten houdt. Wanneer een van deze kritieke services op een beheerde laptop wordt stopgezet via een lokaal administratoraccount, stuurt mijn service direct een melding via e-mail naar het IT- en securityteam. Deze e-mail vermeldt de naam van de laptop, de naam van de service en de datum en het tijdstip van het incident.
De nieuwe service wordt op zijn beurt beschermd door een tweede service, die controleert of de eerste service actief blijft. Het project is geslaagd wanneer het IT-team binnen de verwachte tijd van 2 à 3 minuten een waarschuwingsmail ontvangt nadat een service is stopgezet.


%---------- 3. Technologieën --------------------------------------------------

\section{Technologieën}%
\label{vst:technologieen}

% Met welke talen, frameworks, databanken en diensten ga je werken? Welke
% liggen vast door de omgeving van het bedrijf, en welke kies je zelf? Vermeld
% ook wat nieuw voor je is: technologische verbreding wordt gewaardeerd.
%
% Vermeld hier ook of je AI denkt in te zetten, en waarvoor. Dat mag en wordt
% zelfs verwacht; in je rapport verantwoord je later wat je precies gedaan hebt.
%
% Verwerkt je toepassing persoonsgegevens, of zit er AI in wat je bouwt? Zeg
% dat dan nu al in één zin: dan weet je begeleider meteen dat de AVG of de
% AI Act meespeelt, en kan dat mee bekeken worden bij de goedkeuring.

De gebruikte taal is C sharp, en de gebruikte technologie is Worker Service : Dit is een speciaal projecttype in .NET dat is ontworpen voor langlopende achtergrondtaken. Het maakt gebruik van de klasse BackgroundService en de methode ExecuteAsync. Ik heb dit idee zelf gekozen omdat ik op de hoogte ben van de noden van IT team. Ik werk zelf daarbij. Die technologie is volledig nieuwe voor me. 
Ik ga AI gebruiken om mijn code te verfijnen.
De service gaat geen persoonsgegevens verwerken, die gaat enkel de naam van de gestopte service en toestel vermelden.   

%---------- 4. Afbakening -----------------------------------------------------

\section{Afbakening}%
\label{vst:afbakening}

% DE BELANGRIJKSTE RUBRIEK. Je stage en je graduaatsproject lopen bij hetzelfde
% bedrijf; dit is waar je ze uit elkaar houdt.
%
% Beschrijf:
%   - wat je dagelijkse stagewerk is;
%   - wat het afgebakende geheel is dat JIJ zelfstandig analyseert, ontwerpt
%     en oplevert - dat is je graduaatsproject;
%   - wat er uitdrukkelijk buiten valt.
%
% Toets voor jezelf: welke analyse en welk ontwerp maak jij zelf? Kan je daar
% niets over schrijven, dan heb je een taak en geen project.

Tijdens mijn stage bouw ik een webapplicatie die de onboarding- en offboardingtaken automatiseert. De applicatie communiceert hiervoor met de Microsoft Graph API en Exchange Online. Het doel is dat het IT-team tijd bespaart door deze herhalende taken te automatiseren. De backend wordt ontwikkeld in C sharp en voor de frontend wordt het React-framework gebruikt. Mijn graduaatsproject staat los van de stageopdracht, al maak ik wel gebruik van de bronnen op de werkplek een geregistreerde app in Azure die over Mail.Send-machtigingen beschikt).
Ik heb zowel de analyse als het ontwerp van het project volledig zelfstandig uitgewerkt, op basis van mijn ervaring binnen het IT-team.


%---------- 5. Verwachte meerwaarde -------------------------------------------

\section{Verwachte meerwaarde}%
\label{vst:meerwaarde}

% Wat heeft het bedrijf hieraan? Wat blijft er van jouw werk over wanneer je
% stage afgelopen is? Denk aan tijdwinst, minder fouten, nieuwe mogelijkheden,
% lagere kosten, of een basis waarop verder gebouwd kan worden.

De service is volledig eigendom van het bedrijf en kan eenvoudig worden uitgerold naar alle beheerde toestellen. Zodra een kritieke service wordt stopgezet, waarschuwt de service direct het IT-team, zodat zij onmiddellijk op de hoogte zijn van het incident. Bovendien vormt deze oplossing een schaalbare basis voor toekomstige uitbreidingen, zoals het monitoren van andere cruciale componenten en services, waaronder BitLocker (Disk encryptie op Windows toestellen ).

%---------- 6. Planning en werkbelasting --------------------------------------

\section{Planning en werkbelasting}%
\label{vst:planning}

% Het graduaatsproject is ongeveer tien werkdagen, samen ongeveer 100 uren.
% Verdeel die over de fasen van je project en geef per fase aan wat je oplevert.
% Vermeld ook of je die dagen gefocust opneemt of gespreid over de stageperiode,
% en dat je dit met je mentor afgesproken hebt.

\begin{table}[htbp]
  \centering
  \begin{tabular}{llr}
    \toprule
    \textbf{Fase} & \textbf{Wat lever je op} & \textbf{Uren} \\
    \midrule
    Analyse en vereisten      & overzicht van de vereisten   & 14  \\
    Technologiekeuze          & verantwoorde keuze           & 11  \\
    Ontwerp                   & architectuur en gegevensmodel& 15  \\
    Realisatie                & werkende oplossing           & 40  \\
    Testen en validatie       & testverslag                  & 10  \\
    Rapport, poster en demo   & de op te leveren documenten  & 10  \\
    \midrule
    \textbf{Totaal}           &                              & \textbf{100} \\
    \bottomrule
  \end{tabular}
  \caption[Planning en werkbelasting]{\label{tab:vst-planning}%
    }
\end{table}

%---------- 7. Risico's -------------------------------------------------------

\section{Risico's}%
\label{vst:risicos}

% Wat kan dit project doen mislukken, en wat doe je dan? Denk aan: toegang tot
% gegevens of systemen die je nodig hebt, afhankelijkheid van een collega of
% een derde partij, technologie die je nog moet leren, of drukte op de
% werkvloer waardoor je dagen niet vrijkomen.
%
% Twee tot drie risico's met telkens één zin over wat je terugvalplan is.

1. Een lokale administrator zet de service stop: De service is zo geconfigureerd dat deze in dat geval standaard automatisch opnieuw opstart.
2. Een lokale administrator schakelt de service volledig uit: Een tweede achtergrondservice monitort continu of de eerste service actief is en stuurt direct een     waarschuwing naar het IT-team zodra deze wordt uitgeschakeld.
3. De melding komt niet aan in het ticketingsysteem van het IT-team: Om er zeker van te zijn dat niemand het incident mist, wordt de waarschuwingsmail niet alleen naar het ticketingsysteem gestuurd, maar ook naar een distributielijst waar alle leden van het IT-team deel van uitmaken.



\printbibliography[heading=bibintoc]

\end{document}
